\chapter{The Audio System}
\label{ch:audio-system}

CNA's \cnans{Microsoft::Xna::Framework::Audio} namespace sits on a foundation the project
debated openly before choosing: real XNA runs XACT content through Microsoft's own XACT
engine; FNA reimplements that engine as FAudio. CNA does neither --- the project owner
explicitly rejected an optional-FAudio-backend proposal, and instead built a hand-written
XACT container parser (\texttt{CNA::Internal::Audio::XactParser}) that feeds real,
SDL3\_mixer-based playback. This single decision is the root cause of essentially every
deviation from FNA's own audio behavior catalogued in this chapter.

\section{Why SDL3\_mixer, and what that choice costs}

The CNA mixer uses a fixed reference format: its device mixer always requests 16-bit signed
stereo at 44100~Hz, regardless of what the physical output device's native format actually
is, then queries and logs the format SDL3\_mixer actually negotiated. This is CNA's deliberate
simplification, not a universal SDL3\_mixer restriction. The cost of building on SDL3\_mixer rather than a dedicated 3D-audio engine
shows up consistently: no true per-source 3D audio graph, no auxiliary-send/reverb bus, and a
single shared ``cooked callback'' slot per track used for both filtering and pan-crossfeed at
once, rather than independent processing stages.

The library's object model has three levels. One process-wide \texttt{MIX\_Mixer} owns the
device-side mix; decoded \texttt{MIX\_Audio} objects hold playable audio; and each playback is a
\texttt{MIX\_Track} carrying gain, frequency ratio, loop state, and the cooked callback. A
\cnaclass{SoundEffect} retains decoded audio, while a \cnaclass{SoundEffectInstance} creates and
controls a track.

Mixer creation and lookup are serialized by one mutex. The first use also pins an extra
\texttt{SDL\_INIT\_AUDIO} reference. Without that pin, destroying an SDL3\_mixer device could
drop SDL's global audio-subsystem reference to zero while a
\cnaclass{DynamicSoundEffectInstance} still owned an \texttt{SDL\_AudioStream}; a real ASan
trace ended in \texttt{SDL\_UnbindAudioStream}. \texttt{DestroyMixer()} is implemented and a
generation counter protects stale tracks if it is ever used, but no production caller invokes
it. In the current lifecycle the mixer and its initialization reference live until process
exit.

\section{SoundEffect and its instances}

\cnaclass{SoundEffect} loads decoded file/stream assets or raw PCM and is move-only, specifically so that
disposing the parent cascades correctly to every \cnaclass{SoundEffectInstance} it ever
handed out via \texttt{CreateInstance()}. Its raw-buffer constructors require
\emph{headerless} 16-bit PCM data --- passing an entire WAV file's bytes, RIFF header
included, is flagged directly in the class's own documentation as a common, specific mistake
to avoid. \cnaclass{SoundEffectInstance} exposes volume/pan/pitch/looping/state and a
\texttt{Apply3D(listener, emitter)} method implementing a distance-and-pan
\emph{approximation} rather than full 3D audio --- SDL3\_mixer has no native Doppler
computation, so CNA computes the F3DAudio-compatible Doppler frequency ratio itself and applies
it to the track rather than spatializing samples through a native 3D engine. Pitch shifting uses XNA's exact exponential formula,
$2^{\text{pitch}}$, and a real state-variable filter (low-pass, high-pass, or band-pass) is
implemented via an SDL3\_mixer per-track callback matching FAudio's own filter algorithm
precisely. \cnaclass{DynamicSoundEffectInstance} supports both integer and, as a
\texttt{CNAEXT} addition with no real-XNA equivalent, floating-point buffer submission; its
\texttt{PendingBufferCount} only decreases once a submitted chunk is genuinely consumed by
playback, matching FNA's own behavior rather than decrementing eagerly at submission time.

\subsection{Full method reference: SoundEffect}

\cnaclass{SoundEffect} exposes four static, process-wide mixing knobs, read directly from its
header --- \texttt{MasterVolume}, \texttt{DistanceScale}, \texttt{DopplerScale}, and
\texttt{SpeedOfSound} --- all consulted by every live \cnaclass{SoundEffectInstance} at once,
not per-instance settings. Beyond \texttt{CreateInstance()} and the static, fire-and-forget
\texttt{Play(volume, pitch, pan)} convenience helper already named above, three constructor
families exist: a \texttt{CNAEXT} file-path constructor for loading directly from an asset name
(not part of the real XNA~4.0 API, which only ever loads through
\cnaclass{ContentManager}); two raw-buffer constructors (a full-format one taking explicit
channel count/sample rate/loop points, and a simpler one assuming mono 16-bit PCM at a given
sample rate) --- both requiring \emph{headerless} PCM data as already emphasized;
\texttt{GetSampleDuration} / \texttt{GetSampleSizeInBytes} (both static) convert between a raw
buffer's byte size and its playback duration, in either direction, for PCM data at a given
format --- the calculation a streaming or dynamically-generated sound needs to know how large
a buffer to allocate for a target duration, or how long a given buffer will actually play for.

File and stream construction delegates to \texttt{MIX\_LoadAudio} or
\texttt{MIX\_LoadAudio\_IO} with predecode enabled. CNA's vendored build consequently decodes
WAVE, AIFF, VOC, AU, FLAC through dr\_flac, MP3 through dr\_mp3, Ogg Vorbis through
stb\_vorbis, and MIDI through Timidity. It explicitly disables GME, XMP/MOD, mpg123,
FluidSynth, Opus, WavPack, libFLAC, vorbisfile, and Tremor. In particular, advertised or
third-party \texttt{.opus}, AAC, and WMA files are not made decodable merely by this constructor.
The raw-buffer overloads remain signed 16-bit PCM only.

\begin{lstlisting}
// Worked example: loading a sound effect from raw, headerless PCM data
// (e.g. procedurally synthesized, or decoded by application code from a
// format SoundEffect itself doesn't parse) and playing it once.
std::vector<SharpRuntime::bytecs> pcmData = SynthesizeBeep(440.0f, 0.2f);  // 440 Hz, 200 ms

const SoundEffect beep(pcmData, 44100, AudioChannels::Mono);
beep.Play(1.0f, 0.0f, 0.0f);   // volume, pitch (semitone-equivalent), pan
\end{lstlisting}

\subsection{Full method reference: SoundEffectInstance and 3D audio}

Beyond \texttt{Volume} / \texttt{Pan} / \texttt{Pitch} / \texttt{IsLooped} / \texttt{State} (a
\cnaclass{SoundState}: \texttt{Playing} / \texttt{Paused} / \texttt{Stopped}) and
\texttt{Play} / \texttt{Pause} / \texttt{Resume} / \texttt{Stop} (with an \texttt{immediate} overload
distinguishing a hard cut from a natural fade-out), \texttt{Apply3D(listener, emitter)} is
where this class's real complexity lives. \cnaclass{AudioListener} and \cnaclass{AudioEmitter}
both carry the same four \cnaclass{Vector3} fields --- \texttt{Position}, \texttt{Forward},
\texttt{Up}, \texttt{Velocity} --- mirroring a camera's own transform properties from
Chapter~\ref{ch:graphicsdevice}, plus \cnaclass{AudioEmitter}'s own \texttt{DopplerScale}. Calling \texttt{Apply3D}
even once permanently changes how the instance behaves from then on: reading
\texttt{SoundEffectInstance.cpp} directly confirms it latches an internal \texttt{is3D} flag
that never resets, matching FNA's own equivalent latch --- once \texttt{Apply3D} has run, its
derived attenuation/pan/Doppler values keep governing the actual audio output on every later
\texttt{Play()} / \texttt{setVolumeProperty()}/etc.\ call, not just the call that invoked
\texttt{Apply3D} itself, even though those setters still update their own properties
independently (so a caller reading \texttt{Volume} back still sees what it last set, even
though the real output is governed by \texttt{Apply3D}'s computation instead).

\begin{lstlisting}
// Worked example: positioning a 3D sound effect relative to a moving
// listener -- Apply3D must be called at least once for spatial
// attenuation/pan to take effect at all, and every frame afterward to
// track a moving emitter or listener.
AudioListener listener;
listener.setPositionProperty(cameraPosition);
listener.setForwardProperty(cameraForward);
listener.setUpProperty(Vector3::Up);
listener.setVelocityProperty(cameraVelocity);   // needed for Doppler

AudioEmitter emitter;
emitter.setPositionProperty(explosionWorldPosition);
emitter.setVelocityProperty(Vector3::Zero);      // a stationary emitter

explosionInstance.Apply3D(listener, emitter);
explosionInstance.Play();
\end{lstlisting}

The exact scalar rules are stronger than ``some distance and pan approximation''. With
$d_n=d/\texttt{DistanceScale}$, attenuation is

\[
  a(d_n)=\begin{cases}1,&d_n<1,\\
  \operatorname{clamp}(1/d_n,0,1),&d_n\geq1.
  \end{cases}
\]

This is FAudio's default inverse law: full volume inside \texttt{DistanceScale}, inverse
falloff beyond it. Pan projects emitter displacement onto the listener's actual right vector
$\texttt{Forward}\times\texttt{Up}$, not world~X. For pan $p\in[-1,1]$, the stereo crossfeed
coefficients $(LL,RL,LR,RR)$ are

\[
\begin{cases}
(1+0.5p,-0.5p,0,1+p),&p\leq0,\\
(1-p,0,0.5p,1-0.5p),&p>0.
\end{cases}
\]

Thus a hard pan mixes both source channels into the favoured speaker; it does not merely mute
one input channel. Doppler uses the F3DAudio velocity projections and clamps the resulting
factor to $[0.5,4.0]$, with a NaN guard. The final frequency ratio is
$2^{\text{Pitch}}\times\text{dopplerFactor}\times\texttt{DopplerScale}$, using a default
\texttt{SpeedOfSound} of 343.5. These scalar results persist after \texttt{Apply3D}; later play,
volume, and pitch calls recombine them. What remains approximate is speaker diffusion,
elevation/HRTF, and per-sample spatial filtering, not the documented attenuation and Doppler
formulae themselves.

\subsection{Full method reference: DynamicSoundEffectInstance}

\cnaclass{DynamicSoundEffectInstance} inverts the usual ``load a complete asset, then play
it'' flow into a pull-based streaming model, the correct tool for procedurally generated audio
(synthesized tones, streamed voice/network audio) where the complete sample data does not
exist up front. \texttt{BufferNeeded} (a \texttt{System::EventHandler<System::EventArgs>} ---
Chapter~\ref{ch:sharp-runtime-namespaces} covers this event type in full) fires when the internal queue is running low and
more data should be submitted; \texttt{SubmitBuffer(buffer)} (headerless 16-bit PCM, matching
\cnaclass{SoundEffect}'s own raw-buffer requirement) queues one more chunk, with a second
overload accepting an explicit offset/count sub-range of a larger buffer without copying it
first. The \texttt{CNAEXT} \texttt{SubmitFloatBufferEXT(buffer)} sibling accepts normalized
\texttt{float} samples directly instead, for application code that already generates audio in
floating-point form and would otherwise need to manually quantize to 16-bit integers first ---
real XNA has no equivalent at all, since the strict API only ever accepts integer PCM.
\texttt{PendingBufferCount} (already named above) is the standard way to decide whether
\texttt{BufferNeeded} needs a response right now or the queue still has enough buffered ahead.
The instance owns an \texttt{SDL\_AudioStream} and binds it to a mixer track; the stream is not
owned by the mixer. Playback keeps the track alive while exhausted so future submitted chunks
can resume it. Matching FNA, \texttt{Stop(false)} is not a graceful drain operation here: a
non-immediate stop throws \texttt{InvalidOperationException}; use \texttt{Stop()} for the
supported immediate stop.

\begin{lstlisting}
// Worked example: a procedurally generated tone using the pull-based
// streaming model -- BufferNeeded drives generation, rather than the
// whole sound needing to exist as one complete buffer up front.
DynamicSoundEffectInstance tone(44100, AudioChannels::Mono);
tone.BufferNeeded += [&tone](System::Object*, const System::EventArgs&)
{
    // Keep two chunks queued ahead so playback never starves waiting
    // for the next one to be generated.
    while (tone.getPendingBufferCountProperty() < 2)
    {
        tone.SubmitBuffer(GenerateNextSineChunk(440.0f));
    }
};
tone.Play();
\end{lstlisting}

\subsection{A real, recent bug: a whole chunk popped for one byte of real playback}

The claim earlier in this section --- \texttt{PendingBufferCount} only decreases once a chunk
is \emph{genuinely} consumed --- is worth the same scrutiny as the \cnaclass{Cue} lifetime bug
above, because a real, externally-reported bug (found the same week, independently
re-verified by hand-tracing the exact algorithm against real byte counts before any fix) shows
how easy that guarantee is to get subtly wrong. The pre-fix byte-accounting loop compared
\texttt{total} (bytes SDL reports as consumed so far) against \texttt{queuedBytes} (bytes this
instance still believes are queued), and popped an \emph{entire} submitted chunk off its
internal tracking list the instant \texttt{total > queuedBytes} became true at all --- not once
that specific chunk's own full byte count had actually finished playing. Hand-tracing two
whole-second, 4096-byte chunks against real elapsed time makes the consequence concrete: after
roughly 30~ms of real playback (perhaps 92 bytes genuinely consumed), the comparison already
went true, and the \emph{entire} first 4096-byte chunk was popped --- \texttt{PendingBufferCount}
dropped from~2 to~1 after 30~ms of a chunk that needed a full second to actually finish. By
roughly 1017~ms (barely past that first chunk's own real duration), the loop's second iteration
fired for the same reason, popping the \emph{second} chunk too --- \texttt{PendingBufferCount}
reached~0 a little over a second in, roughly half the real two-second runway both chunks
together should have provided.

The fix computes a single \texttt{consumed = total - queuedBytes} budget once per
\texttt{Update()} call, and only pops a chunk once that budget actually covers the chunk's
entire real size, decrementing the budget as each one is confirmed --- so a second chunk can
never be credited with bytes the first one has not finished consuming yet. The practical
consequence for the worked example above: before this fix, the \texttt{while
(getPendingBufferCountProperty() < 2)} refill loop could be triggered roughly twice as often as
the audio hardware actually needed, silently over-generating and over-submitting chunks well
ahead of genuine playback need --- wasted CPU work rather than an audible glitch, and exactly
the kind of gap invisible to any test that does not model real elapsed time against real chunk
durations, which is why it went untested for as long as it did.

\section{The XACT object model: AudioEngine, SoundBank, WaveBank, Cue}

\cnaclass{AudioEngine} parses a compiled \texttt{.xgs} settings file and owns the lifecycle of
every \cnaclass{SoundBank} / \cnaclass{WaveBank} / \cnaclass{Cue} built from it, including
per-category instance-limit enforcement and both global and cue-scoped XACT runtime
variables. \cnaclass{WaveBank} supports both a fully in-memory mode and a real streaming mode
(reading by offset and packet size rather than loading everything up front).
\cnaclass{SoundBank::Dispose()} force-stops every cue it ever produced, matching the real
FACT engine's own \texttt{FACTSoundBank\_Destroy} behavior exactly. \cnaclass{Cue} implements
the full XACT playback state machine (\texttt{Created}, \texttt{Preparing}, \texttt{Prepared},
\texttt{Playing}, \texttt{Stopping}, \texttt{Stopped}) with an independent \texttt{paused}
flag layered on top of \texttt{Playing} --- matching FACT's own bitmask-shaped state rather
than a naive single enum --- real fade-in/fade-out ramps, per-tick re-evaluation of
Runtime Parameter Control curves, and all five of XACT's randomized track-variation selection
algorithms (\texttt{Ordered}, \texttt{OrderedFromRandom}, \texttt{Random},
\texttt{RandomNoRepeats}, \texttt{Shuffle}).

\subsection{Full method reference and a worked example: playing a cue}

\cnaclass{AudioEngine}'s own public surface beyond construction (from a compiled
\texttt{.xgs} settings file) is small and mostly about the two global runtime-variable
categories XACT content can reference: \texttt{GetGlobalVariable} / \texttt{SetGlobalVariable}
read/write an engine-wide named variable (an RPC curve or a piece of gameplay code can both
read the same value --- a common use is a global ``player health'' or ``distance to danger''
variable that several different sound cues' volume/pitch RPC curves all key off of at once),
\texttt{GetCategory(name)} returns an \cnaclass{AudioCategory} handle for bulk
\texttt{Pause} / \texttt{Resume} / \texttt{Stop} / \texttt{SetVolume} control over every cue in one
named category simultaneously (the standard technique for a game-wide ``mute music'' or
``duck all SFX during a cutscene'' toggle), and \texttt{Update()} must be called once per
frame --- it is what actually sweeps finished fire-and-forget cues and advances every active
cue's fade/RPC state, mirroring FACT's own \texttt{FACTAudioEngine\_DoWork} call.
\cnaclass{SoundBank::GetCue(name)} returns an owned \cnaclass{Cue} the caller controls
explicitly (holds a reference, can call \texttt{SetVariable}, decides when to
\texttt{Play} / \texttt{Stop} it); \texttt{PlayCue(name)} and its 3D overload,
\texttt{PlayCue(name, listener, emitter)}, are the fire-and-forget alternative --- the bank
creates and owns the \cnaclass{Cue} internally, and \texttt{AudioEngine::Update()}'s sweep
(named above) is what eventually cleans it up once playback finishes, not the caller.

\begin{lstlisting}
// Worked example: a footstep cue whose pitch varies with the player's
// current speed, via a global XACT runtime variable an RPC curve
// authored in the XACT project reads -- the standard way gameplay code
// drives audio parameters without the audio system needing to know
// anything about game-specific concepts like "speed" directly.
AudioEngine engine("Content/Audio/GameAudio.xgs");
SoundBank soundBank(&engine, "Content/Audio/Sounds.xsb");
WaveBank waveBank(&engine, "Content/Audio/Sounds.xwb");

// Each frame, before any cue that reads "PlayerSpeed" plays:
engine.SetGlobalVariable("PlayerSpeed", currentPlayerSpeed);

// Fire-and-forget: the bank owns this Cue's lifetime from here on.
soundBank.PlayCue("footstep");

// A held, caller-managed Cue instead, for a looping engine sound that
// needs to be stopped explicitly later:
Cue* engineLoop = soundBank.GetCue("engine_loop");
engineLoop->Play();
// ... later, when the vehicle is destroyed:
engineLoop->Stop(AudioStopOptions::AsAuthored);   // respects authored fade-out
delete engineLoop;

// Once per frame, unconditionally:
engine.Update();
\end{lstlisting}

\subsection{A real, recent bug: a Cue obtained but never played could outlive its bank}

The \texttt{GetCue} / \texttt{Play} split shown above --- obtaining a \cnaclass{Cue} handle
without immediately playing it --- is precisely the case a real, externally-reported bug (dated
2026-07-17, independently re-verified line-by-line before any fix was written, per this
project's own standard practice) found a genuine gap in. \cnaclass{SoundBank}'s own
\texttt{Dispose()} force-stops every cue the bank still owns by walking its
\texttt{activeCues\_} list --- but \cnaclass{Cue}'s constructor never registered a newly
obtained cue into that list at all; only \texttt{Play()} did, at each of its three real exit
points. A cue obtained via \texttt{GetCue()} but never played --- exactly \texttt{engineLoop}
above, in the brief window before its own \texttt{Play()} call runs --- was therefore
completely invisible to the bank's disposal cascade: disposing the \cnaclass{SoundBank} while
such a cue was still held would leave it referencing a \texttt{SoundBank*} that could go
dangling the moment the bank object itself was actually destructed, with no exception and no
warning at either the bank's disposal or the cue's own subsequent (now unsafe) \texttt{Play()}
call.

The fix moves registration into \cnaclass{Cue}'s own constructor, unconditionally, so a cue is
tracked by its bank for its entire C\texttt{++} lifetime --- prepared, playing, stopping, or
stopped --- not merely while actively playing; unregistration moved correspondingly from the
moment playback genuinely stops to \cnaclass{Cue::Dispose()} instead, closing both the
never-played gap and a narrower played-then-stopped-but-not-yet-disposed gap in the same
change. The practical rule this implies for any code holding a caller-managed \cnaclass{Cue}
the way \texttt{engineLoop} does above: the safety net now genuinely covers the full period
between \texttt{GetCue()} and your own \texttt{delete}, not merely the period the cue happens
to be audibly playing --- but only as of this fix; a build predating it would leave the same
\texttt{engineLoop} pattern one \texttt{soundBank.Dispose()} away from a dangling pointer if
the vehicle were destroyed before the cue's own \texttt{Play()} call had run.

\subsection{AudioCategory: a detail easy to get wrong}

\cnaclass{AudioCategory}'s own class-level doc comment states a precise, checkable behavior
worth calling out on its own, since it is easy to assume otherwise: \texttt{SetVolume(volume)}
does not merely set a baseline for \emph{future} \texttt{Play()} calls in that category --- it
retroactively re-applies to every cue in the category that is already playing, immediately.
A gameplay system that fades a category's volume down over several frames (rather than
snapping it in one call) genuinely audibly attenuates whatever is already playing, frame by
frame, not just whatever starts playing next.

\begin{lstlisting}
// Worked example: ducking music volume while a cutscene dialogue cue
// plays, then restoring it -- SetVolume's retroactive behavior is
// exactly what makes a smooth per-frame fade of already-playing music
// possible, rather than only affecting cues that start after the call.
AudioCategory music = engine.GetCategory("Music");
music.SetVolume(duckedVolume);   // audibly fades already-playing music now

dialogueCue->Play();
// ... later, once the dialogue cue finishes:
music.SetVolume(1.0f);           // restores full volume, again retroactively
\end{lstlisting}

\cnaclass{AudioEngine::RendererDetails} (a \texttt{std::vector<RendererDetail>}, each exposing
a \texttt{FriendlyName} / \texttt{RendererId} pair) exists to enumerate available audio output
renderers, mirroring real XNA's own multi-renderer-aware API shape --- but reading
\texttt{RendererDetail.hpp}'s own source comment directly shows this is currently a shape
without substance in CNA specifically: ``production code only ever constructs the single
hardcoded SDL3\_mixer renderer'' (a second, distinct instance exists only in test code, to
exercise \texttt{Equals} / \texttt{operator==}'s not-equal path at all). Application code that
enumerates \texttt{RendererDetails} expecting a real per-device choice, the way a Windows XAudio2
game might see one entry per real audio device, will find exactly one entry, always the same
one, regardless of what audio hardware is actually attached.

\subsection{SoundEffect::FromStream, and the exception types this namespace defines}

Beyond the raw-buffer and asset-name constructors already covered,
\texttt{SoundEffect::FromStream(stream)} is a fourth, static loading path: it parses a
complete WAV \emph{container} (RIFF header included) directly from a
\texttt{std::istream}, returning an owning pointer the caller is responsible for deleting ---
the one \cnaclass{SoundEffect} construction path that \emph{wants} the container header,
in explicit contrast to every other constructor's headerless-PCM requirement named earlier in
this section.

This namespace defines three audio-specific exception types, but only two represent active
throw paths. \cnaclass{InstancePlayLimitException} is declared and its constructors are tested,
yet no production implementation throws it. Fire-and-forget \texttt{SoundEffect::Play} creates
a fresh \texttt{MIX\_Track} per call and has no simultaneous-instance cap.

\cnaclass{NoAudioHardwareException} and \cnaclass{NoMicrophoneConnectedException} are thrown
when no playback or capture device is available, respectively; both wrap a plain error message
and an optional inner exception, matching the standard .NET exception-chaining constructor
shape --- but not from the same base class. Two of the three,
\cnaclass{InstancePlayLimitException} and \cnaclass{NoAudioHardwareException}, derive from a
COM-interop-style external exception base, matching real XNA's own choice to model those two
conditions as externally-originating errors.

\cnaclass{NoMicrophoneConnectedException}, the third, derives from a plain exception base
instead --- a real, if minor, asymmetry across three otherwise structurally similar exception
types. Its existence does not make \cnaclass{InstancePlayLimitException} operational.

\section{A real user-reported bug, and the audit it triggered}

The most recent, and most consequential, audit of this namespace was not routine --- it was
triggered by real user reports of high-pitched, sped-up, distorted, or missing audio in an
actual ported game. The highest-priority finding directly explains the ``missing audio''
class of report: the XNB \cnaclass{SoundEffectReader} only ever accepted a narrow PCM format,
meaning \emph{any} MS-ADPCM-compressed XACT wave bank or XNB sound effect --- a common,
space-saving compression choice for XACT content --- silently failed to decode at all,
because the reader's synthetic WAV wrapper was missing the coefficient table MS-ADPCM
decoding requires. This is now fixed, and the fix was extended to cover 8-bit PCM, IEEE
float, and IMA-ADPCM as well, not just the one format that triggered the investigation. The
current XNB reader's exact set is therefore PCM 16-bit (direct), PCM 8-bit, IEEE float 32-bit,
MS-ADPCM 4-bit, and IMA-ADPCM 4-bit (the latter four wrapped for SDL decoding), restricted to
mono or stereo. XMA2 remains explicitly rejected. For XNA-authored MS-ADPCM with
\texttt{cbSize=0}, the wrapper synthesizes the standard seven coefficient pairs before decode.

A second real fix addressed an asymmetry in how \cnaclass{DynamicSoundEffectInstance}
handled integer versus float buffer submission --- a genuine code-level risk, not merely a
missing feature. And because no end-to-end audible-output conformance test existed at all
before this audit, a new deterministic offline rendering harness was built specifically to
measure this class of bug directly: it empirically rules out SDL3\_mixer's own resampler as
the cause of the reported pitch issue, confirming frequency is preserved within 0.1\% across
22050/44100/48000/96000~Hz source material resampled into the fixed 44100~Hz mixer, provided
the source sample rate is declared correctly.

A separate, real stereo-panning bug used to eliminate the opposite source channel rather than
crossfeed both inputs into the favoured speaker. The current instance \texttt{Play},
\texttt{Apply3D}, pan setter, and static fire-and-forget
\texttt{SoundEffect::Play(volume,pitch,pan)} paths all use the cooked-callback crossfeed model;
the convenience helper allocates a small per-track pan state and defers freeing it until the
stop callback has safely returned.

A follow-up round of the same audit, prompted by nothing external this time --- simply
re-reading \cnaclass{SoundEffect} / \cnaclass{SoundEffectInstance} /
\cnaclass{DynamicSoundEffectInstance} end to end against FNA once more --- found a third,
wide-reaching numeric bug in exactly the \texttt{Pitch} formula this chapter's own opening
section states as correct today: \texttt{Pitch}'s conversion to a playback-rate ratio was, for a
real stretch of this project's history, \texttt{ratio = (pitch<0) ? (1+pitch*0.5) :
(1+pitch)} --- a linear approximation, not real FNA's genuine exponential octave curve,
$2^{\text{pitch}}$, independently re-verified directly against FNA's own
\texttt{SoundEffectInstance.cs}. The two formulas agree only at the three points
$\text{pitch} \in \{-1, 0, 1\}$; at $\text{pitch} = 0.5$ the old linear formula gave a ratio of
exactly $1.5$ against the correct $2^{0.5} \approx 1.4142$ --- roughly one semitone off, clearly
audible, and never caught by any existing test because every prior Doppler/pitch test happened
to exercise only $\text{pitch} = 0$, the one value (besides $\pm 1$) where the two formulas
coincidentally produce the same answer. Because \cnaclass{Cue} routes every XACT pitch
application --- base pitch, an RPC pitch curve, effect-variation pitch --- through this same
\texttt{SoundEffectInstance::setPitchProperty()} internally, the bug reached all pitched XACT
playback, not only code calling \cnaclass{SoundEffectInstance.Pitch} directly. The fix unified
what had been three separately duplicated copies of the wrong formula (the static
\texttt{SoundEffect::Play} fire-and-forget path, the shared \texttt{ApplyTrackProperties()}
helper both \texttt{Play()} and \texttt{Apply3D()} use, and the pitch setter itself) into one
small, independently-testable, correctly-exponential helper --- and, while fixing it, caught a
second small defect in the same code: \texttt{ApplyTrackProperties()}'s own doc comment already
\emph{claimed} to implement the correct $2^{\text{pitch}}$ formula before the fix, a stale,
aspirational comment describing behavior the code had never actually had.

\section{Accepted deviations, stated as such}

Several differences from a hypothetical ``full XACT'' implementation are documented as
partial or accepted behavior. The \texttt{REPLACE\_QUIETEST} instance-limit branch is
unfinished, while queue and replace-oldest share behavior inherited from the FACT/FNA model;
do not present the three policies as independently complete. A cue-level instance limit has no category- or
same-cue victim-search filter. Runtime Parameter Control curves targeting a DSP preset are
never evaluated, because no DSP preset system exists at all. \texttt{Apply3D}'s panning is a
single-axis linear approximation, not full multi-speaker 3D diffusion. A parsed per-track
filter can only ever decode as low-pass or high-pass, never band-pass --- itself a real bit-decoding
quirk in FAudio being faithfully replicated, not invented independently. Reverb is a
documented no-op, since SDL3\_mixer provides no auxiliary send/return bus to implement it
against, and there is no HRTF or elevation modeling in the 3D audio approximation at all.

\section{Real memory-safety bugs found and fixed}

Beyond behavioral correctness, the same audit pass found and fixed several genuine
memory-safety issues through a systematic ``what happens if operation X runs concurrently
with operation Y'' review: a mixer generation-counter pair with a real use-after-free,
a heap-buffer-overflow in XACT name parsing, and a genuine data race in \cnaclass{WaveBank}'s
own cache --- \texttt{Dispose()} could race a concurrently in-flight
\texttt{GetSoundEffect()} decode --- fixed with a mutex. Two crash findings remain open at the
pinned revision, but they are not the same shape.  \texttt{AUD-15-021} is an intermittent
late failure after a long audio-filter run; \texttt{AUD-15-022} is deterministic only for a
narrow cross-suite Cue/DynamicSoundEffectInstance ordering at process teardown.  ASan and TSan
perturb both reproductions, while \texttt{gdb} and Valgrind were unavailable for the recorded
campaign.  Neither finding is evidence that every ordinary full-suite run crashes.

\section{Testing what reaches the mixer}

\texttt{OfflineAudioRenderer} is a test helper with unusually broad value: it creates a
device-less \texttt{MIX\_Mixer} and calls \texttt{MIX\_Generate} to render real decode,
resample, track properties, callbacks, and mixing into memory. It needs neither an audio device
nor \texttt{SDL\_AUDIODRIVER}, and advances by requested sample count rather than wall-clock
time. Tests can therefore measure pitch, channel energy, duration, and waveform output
deterministically on a headless machine.

This is stronger than asserting which setter was called and narrower than perceptual audio
quality. It proved, for example, that SDL3\_mixer preserves frequency within the chosen
tolerance when 22050, 44100, 48000, and 96000~Hz sources are correctly declared and resampled
to the fixed 44100~Hz mix. That evidence ruled out the resampler as the reported high-pitch
cause and redirected the audit to the incorrect linear pitch formula.

\section{Microphone}

\cnaclass{Microphone} provides real SDL3 capture-device enumeration and streaming, a
\texttt{BufferReady} event, and enforces a buffer-duration constraint of 100 to 1000
milliseconds in multiples of 10 milliseconds --- matching real XNA's own documented
constraint on this API exactly.

\subsection{Full method reference and a worked example}

Device discovery is entirely static: \texttt{Microphone::All} (a
\texttt{std::vector<Microphone*>}) enumerates every capture device SDL3 reports,
\texttt{Microphone::Default} returns the platform's default one (or \texttt{nullptr} if none
exists), and neither list ever needs to be freed by the caller --- every \cnaclass{Microphone}
instance is owned internally and lives for the process's lifetime. Per-instance,
\texttt{State} (a \cnaclass{MicrophoneState}: \texttt{Started} / \texttt{Stopped}) and
\texttt{IsHeadset} / \texttt{SampleRate} describe the device; \texttt{BufferDuration} is the
get/set threshold controlling how often \texttt{BufferReady} fires --- reading its setter's own
doc comment confirms the 100--1000ms-in-multiples-of-10 constraint is enforced by actually
throwing \texttt{System::ArgumentOutOfRangeException} for an out-of-range value, not merely
documented and left unchecked. \texttt{Start()} / \texttt{Stop()} begin/end capture;
\texttt{GetData(buffer)} (and its offset/count sub-range overload) drains whatever captured
audio is currently queued into a caller-supplied buffer, returning the number of bytes
actually written --- the correct way to read \texttt{BufferReady}'s payload, since the event
itself carries no data, only a notification that some is now available.
\texttt{GetSampleDuration} / \texttt{GetSampleSizeInBytes} mirror \cnaclass{SoundEffect}'s own
pair, converting between a raw capture buffer's byte size and its duration for this specific
device's sample rate.

\begin{lstlisting}
// Worked example: capturing microphone input and feeding it straight
// into a DynamicSoundEffectInstance for live playback -- the standard
// "hear yourself" or voice-chat-passthrough pattern, bridging capture
// and playback through the same headerless-PCM buffer format both
// classes share.
Microphone* mic = Microphone::getDefaultProperty();
DynamicSoundEffectInstance playback(mic->getSampleRateProperty(), AudioChannels::Mono);

mic->BufferReady += [mic, &playback](System::Object*, const System::EventArgs&)
{
    std::vector<SharpRuntime::bytecs> captured(4096);
    const int bytesRead = mic->GetData(captured);
    if (bytesRead > 0)
    {
        captured.resize(bytesRead);
        playback.SubmitBuffer(captured);
    }
};

mic->Start();
playback.Play();
\end{lstlisting}

\section{Status}

The audio surface is substantial: real SDL3\_mixer playback, deterministic offline mixing,
streaming dynamic audio, capture, bounded XACT parsing, wave banks, sound banks, cues,
categories, variables, filters, and partial instance-limit behavior. It is not a closed
conformance claim. \texttt{plan\_audio.md}'s three user-reported release blockers---high-pitch
effects, missing/failing loads, and distorted audio---remain unchecked at the pinned revision.
The corrected format readers and pitch/pan fixes materially narrow those reports, but the plan
itself has not declared them resolved. Reverb/aux send is a permanent no-op,
\texttt{REPLACE\_QUIETEST} is partial, \cnaclass{InstancePlayLimitException} has no throw site,
and production mixer teardown is unreachable.
